% GENERATED by evaluation/mini-src/prompts_to_tex.py from approach/src/llm_sad_sam/linkers/experimental/s_linker126.py.
% Do not edit by hand: rerun prompts_to_tex.py, then sync_paper.py.
\section{Prompt Templates}
\label{app:prompts}

\approach uses five \ac{LLM} prompts. We give each prompt verbatim as the implementation sends it. Data that differs per call is shown in braces.

\subsection{Alias Discovery}

\paragraph{Alias extraction.}
\begin{quote}\small\ttfamily\raggedright
Find all alternative names used for these components in the document.\\[\medskipamount]
COMPONENTS: {\normalfont\itshape\{component names, comma-separated\}}\\[\medskipamount]
Find surface forms the document uses to refer to a single named component (introduced short forms, alternate names, or words of multi-word names when they alone clearly mean the full name). Reject terms whose ordinary English use dominates.\\[\medskipamount]
A fragment of a longer identifier is not an alias: if a term appears only as part of a compound or qualified name, do not include it.\\[\medskipamount]
DOCUMENT:\\
{\normalfont\itshape\{document sentences, one per line\}}\\[\medskipamount]
Return JSON:\\
\{\\
\hspace*{1.0em}"abbreviations": [\{"term": "short\_form", "component": "FullComponent"\}],\\
\hspace*{1.0em}"synonyms":      [\{"term": "specific\_alternative\_name", "component": "FullComponent"\}]\\
\}\\
JSON only:
\end{quote}

\paragraph{Alias judge.}
\begin{quote}\small\ttfamily\raggedright
JUDGE: Review these component name mappings for correctness.\\[\medskipamount]
COMPONENTS: {\normalfont\itshape\{component names, comma-separated\}}\\[\medskipamount]
PROPOSED MAPPINGS:\\
{\normalfont\itshape\{proposed mappings as JSON\}}\\[\medskipamount]
An alias is valid when the document establishes an equivalence between a phrase and a single named component. It is invalid when the phrase is generic vocabulary or identifies anything other than that one component. When uncertain, prefer APPROVE.\\[\medskipamount]
Return JSON, echoing each approved mapping in full:\\
\{"approved": [\{"term": "term1", "component": "FullComponent"\}]\}\\
JSON only:
\end{quote}

\subsection{\routeOne}

\paragraph{\judgeOne (name judge).}
\begin{quote}\small\ttfamily\raggedright
Validate components in a document.\\[\medskipamount]
COMPONENTS: {\normalfont\itshape\{component names, comma-separated\}}\\[\medskipamount]
A trace link holds between a sentence and a component when the sentence makes an architectural claim about that component -{}- when it says something about that component as a participant in the system this document describes. Referring to the component as a participant is itself such a claim, even when the sentence says nothing further about it.\\[\medskipamount]
Every case gives you the expression the sentence uses, the sentence itself, the evidence the document supplies, and the component whose name that expression reaches. The evidence says what the expression is doing here; none of it is a verdict.\\[\medskipamount]
\hspace*{1.0em}written -{}- how the sentence writes the component's name. Use exact when the component's full catalog name is written as a name; alias when the document established an alternate form for that component and this sentence writes it; part when only one word of the component's multi-word name is written; and qualified name when the full catalog name occurs only inside a longer joined or dotted identifier. A shorter surface leaves more readings open; it does not make the reading in front of you wrong. Where the sentence does not write the name as such, ask what the expression itself denotes in its local context: a participant in the system, or something merely associated with software.\\[\medskipamount]
Reject only on a positive ground -{}- that the sentence asserts nothing of this component, because the name is doing some other job here, or because the sentence denies what it would otherwise say of it.\\[\medskipamount]
Some sentences use an ordinary English word that happens to coincide with a component's name. Approve only when the sentence uses that word as the name of the component; if it is used in its ordinary sense and the component is not what the sentence is talking about, reject. Capitalization is evidence for a name and its absence is evidence against, but neither settles it on its own.\\[\medskipamount]
An expression that occurs only as part of a longer joined or dotted identifier is naming a piece of that identifier, not a participant in what the sentence describes. An expression denoting what a component acts on or produces refers to that thing and not to the component, however clearly the component is the one acting on it.\\
{\normalfont\itshape\{SENTENCES table of the nearby sentences, included when a case writes only one word of a name\}}\\
Where the sentence does not write the name in full, that a surface can name this component is not evidence that it does here.\\[\medskipamount]
For each case, first quote the EXACT words from the sentence the verdict rests on -{}- the words that state the architectural claim about the component, or "none" if the sentence makes no such claim -{}- then decide approve true/false based on that quote.\\[\medskipamount]
CASES:\\
Case {\normalfont\itshape\{i\}}: "{\normalfont\itshape\{span\}}" -\textgreater{} {\normalfont\itshape\{component name\}}\\
\hspace*{1.0em}{\normalfont\itshape\{[prev: previous sentence] \}}"{\normalfont\itshape\{sentence\}}"\\
\hspace*{1.0em}Evidence: written={\normalfont\itshape\{exact | alias | part | qualified name\}}\\
{\normalfont\itshape\{further cases\}}\\[\medskipamount]
Return JSON:\\
\{"validations": [\{"case": 1, "claim": "\textless{}exact quote or none\textgreater{}", "approve": true\}]\}\\
JSON only:
\end{quote}

\subsection{\routeTwo}

\paragraph{Coreference resolver.}
\begin{quote}\small\ttfamily\raggedright
Resolve references (pronouns and noun phrases that refer back) to components.\\[\medskipamount]
COMPONENTS: {\normalfont\itshape\{component names, comma-separated\}}\\[\medskipamount]
SENTENCES (the document text the cases are drawn from)\\
{\normalfont\itshape\{document sentences as JSON\}}\\[\medskipamount]
For each TARGET sentence below, identify any pronoun or noun phrase in THAT sentence\\
that refers back to a component listed above. Read the TARGET's context in SENTENCES.\\
If a target sentence has no such reference to a listed component, return no resolution\\
for it. Be conservative \textemdash{} only include resolutions you are CERTAIN about.\\[\medskipamount]
Quote the referring expression first, then name the component it points to.\\[\medskipamount]
-{}-{}- Case {\normalfont\itshape\{i\}} -{}-{}-\\
TARGET S{\normalfont\itshape\{n\}}: {\normalfont\itshape\{sentence\}}\\
{\normalfont\itshape\{further cases\}}\\[\medskipamount]
Resolve when the surrounding sentences make one component the clear antecedent, under any form the document uses for it. Avoid resolving when two or more equally plausible antecedents exist.\\[\medskipamount]
Return JSON:\\
\{"resolutions": [\{"case": 1, "sentence": N\_INTEGER, "reference": "the server", "candidates": ["Name", "OtherName"], "component": "Name", "antecedent\_sentence": M\_INTEGER, "antecedent\_text": "exact quote with component name"\}]\}\\[\medskipamount]
JSON only:
\end{quote}

\paragraph{\judgeTwo (coreference judge).}
\begin{quote}\small\ttfamily\raggedright
Validate components in a document. Check coref resolution: does the referring expression in this sentence actually refer to the named component as an architectural participant?\\[\medskipamount]
COMPONENTS: {\normalfont\itshape\{component names, comma-separated\}}\\[\medskipamount]
These are coreference links: a pronoun or noun phrase in the sentence is claimed to refer back to the component, which is NOT named in the sentence itself. Approve only when the sentence contains a genuine referring expression that unambiguously points to THIS component and makes an architectural claim about it. Reject when there is no such referring expression or when the antecedent could equally be a different component. An expression denoting what a component acts on or produces refers to that thing and not to the component, however clearly the component is the one acting on it. When uncertain, reject.\\[\medskipamount]
For each case, first quote the EXACT words from the sentence that state the\\
architectural claim about the component (or write "none" if the sentence makes no\\
such claim), then state the strongest ground there is for rejecting this case under the\\
rules above (or "none" if there is none), then decide: approve unless that ground is one\\
the rules above make decisive. An objection you could raise against most sentences is not\\
a ground for rejecting this one.\\[\medskipamount]
CASES:\\
Case {\normalfont\itshape\{i\}}: pronoun/role-ref -\textgreater{} {\normalfont\itshape\{component name\}}\\
\hspace*{1.0em}Claimed reference: "{\normalfont\itshape\{reference\}}"\\
\hspace*{1.0em}Claimed antecedent (S{\normalfont\itshape\{n\}}): "{\normalfont\itshape\{antecedent text\}}"\\
\hspace*{1.0em}{\normalfont\itshape\{[prev: previous sentence] \}}"{\normalfont\itshape\{sentence\}}"\\
{\normalfont\itshape\{further cases\}}\\[\medskipamount]
Return JSON:\\
\{"validations": [\{"case": 1, "claim": "\textless{}exact quote or none\textgreater{}", "objection": "\textless{}strongest ground to reject, or none\textgreater{}", "approve": true\}]\}\\
JSON only:
\end{quote}
